\documentclass[12pt]{article}
\usepackage[margin=0.8in]{geometry}
\usepackage[english]{babel}
\hyphenpenalty=10000
\exhyphenpenalty=10000
\usepackage[utf8]{inputenc}
\usepackage{setspace}
\usepackage{graphicx}
\usepackage{listings}
\usepackage{xcolor}
\usepackage{fancyhdr}

% Configuracion de codigo
\lstset{
    language=Java,
    basicstyle=\ttfamily\small,
    keywordstyle=\color{blue},
    commentstyle=\color{gray},
    stringstyle=\color{red},
    breaklines=true,
    showstringspaces=false,
    tabsize=2,
    numbers=left,
    numberstyle=\tiny,
    frame=single,
    backgroundcolor=\color{white}
}

% --------------------------------------------------
\begin{document}

\begin{center}
    {\LARGE \textbf{Documentación: Patrones de Diseño}}\\[0.4em]
    {\large \textit{Space Invaders — Proyecto de Programación}}\\[0.3em]
    {\small Fecha: 15 de abril, 2026}\\[0.2em]
    {\small Grupo: Cachopin}\\[0.2em]
    {\small Sprint 2}
\end{center}

\vspace{0.5em}
\hrule
\vspace{1.5em}
\onehalfspacing
\newpage
\tableofcontents
\newpage

\section{Introducción}

En este proyecto se han utilizado tres patrones de diseño fundamentales para que el código sea más mantenible y fácil de extender. Como cualquier proyecto grande, fue necesario tomar decisiones de arquitectura que permitieran agregar nuevas naves, nuevos tipos de disparo y nuevas funcionalidades sin tener que reescribir todo desde cero.

Los tres patrones principales son: \textbf{Factory}, \textbf{Strategy} y \textbf{Composite}. Factory centraliza la creación de naves, Strategy permite tener diferentes tipos de disparos intercambiables, y Composite proporciona una estructura unificada para representar tanto naves como disparos de cualquier complejidad. Cada uno resuelve un problema específico dentro del sistema.

\newpage
\section{Factory: La Fábrica de Naves}

\subsection{La Idea}

En lugar de crear naves esparcidas por todo el código, toda la lógica de creación vive en un único lugar. Así, si un día necesitas cambiar cómo se inicializa una nave, solo cambias en un lugar.

En la práctica esto significa:
\begin{itemize}
    \item No duplicas código de inicialización en varios lugares
    \item Cambios centralizados: si la velocidad inicial cambia, lo cambias una sola vez
    \item Agregar nuevas naves es simple y sin afectar código existente
\end{itemize}

\subsection{Cómo está implementado}

La clase \texttt{NaveFactory} es un Singleton, lo que significa que existe una única instancia durante todo el juego. Todas las naves se crean a través de este punto centralizado, con posiciones y velocidades iniciales definidas aquí.

\subsection{Cómo Funciona en Práctica}

La clase NaveFactory es bastante sencilla. Tiene un atributo privado que guarda la única instancia de la fábrica. Su constructor es privado, lo que impide que se cree directamente con \texttt{new}. La única forma de obtenerla es llamando a un método estático que verifica si ya existe; si no, la crea, y si ya existe, la devuelve directamente.

Tiene otro método que recibe un string indicando qué tipo de nave se quiere. Aquí es donde pasa la magia: dependiendo del nombre recibido, devuelve una nueva instancia de Nave1, Nave2 o Nave3. Todos los parámetros iniciales como la posición (50, 55) y la velocidad (1) están definidos aquí de forma centralizada.

Cuando el juego necesita crear una nave, primero obtiene la fábrica llamando al método estático. La primera vez que se llama, se crea la única instancia que existirá durante toda la ejecución. Las siguientes veces, simplemente devuelve la que ya existe.

Luego, se le pide que genere una nave específica indicando su nombre. La fábrica entonces crea la nave con los parámetros predeterminados, devolviendo un objeto del tipo base Naves. Internally, cada nave configura sus propias estrategias de disparo permitidas y construye su estructura interna de píxeles.

\subsection{Ventajas del Patron Factory}

\begin{enumerate}
    \item \textbf{Encapsulacion:} La logica de creacion queda centralizada. Cambios en parametros iniciales se hacen en un único lugar.
    
    \item \textbf{Bajo acoplamiento:} El cliente solo depende de la interfaz \texttt{Naves}, no de las clases concretas.
    
    \item \textbf{Extensibilidad:} Agregar una nueva nave requiere solo agregar un nuevo tipo y su creación en la fábrica, sin afectar el código existente.
    
    \item \textbf{Singleton integrado:} Garantiza una única instancia de fábrica durante toda la ejecución.
    
    \item \textbf{Control global:} Los parámetros de inicialización se pueden ajustar fácilmente en un único lugar.
\end{enumerate}

\newpage
\section{Strategy: Disparos Intercambiables}

\subsection{La Idea}

Cada nave tiene acceso a diferentes tipos de disparo. Unos tienen munición infinita, otros limitada. Unos son un pixel, otros forman figuras más complejas. Strategy permite que cada nave tenga sus armas propias sin duplicar código.

\subsection{Cómo está implementado}

Todas las estrategias de disparo implementan la misma interfaz \texttt{StrategyDisparo}. Cada una sabe cómo crear su disparo característico y gestionar su propia munición.

\subsection{Implementacion en el Proyecto}

En \textit{Space Invaders}, el patron Strategy se emplea para gestionar los distintos tipos de disparo disponibles. Se define una interfaz común \texttt{StrategyDisparo} que todas las estrategias deben implementar, incluida la creacion del disparo con su forma correspondiente.

Todas las estrategias de disparo implementan la misma interfaz. Cada una sabe cuánta munición tiene, si aún le quedan tiros disponibles, y cómo crear un nuevo disparo en una posición específica.

\subsection{Tres Tipos de Disparos}

\subsubsection{Pixel - El Clásico}

Simple, sin límite. Siempre disponible. Coloca un punto en el espacio.

\textbf{Características:}
\begin{itemize}
    \item Infinita
    \item Un pixel
    \item Siempre funciona
\end{itemize}

\subsubsection{Flecha - El Medium}

30 tiros de munición. Forma de V invertida (3 píxeles). El punto medio entre poder y costo.

\textbf{Características:}
\begin{itemize}
    \item 30 tiros
    \item Forma de V
    \item Se acaba
\end{itemize}

\subsubsection{Rombo - El Pesado}

20 tiros. Forma de rombo (5 píxeles). Potente pero se acaba rápido.

\textbf{Características:}
\begin{itemize}
    \item 20 tiros
    \item Forma de rombo
    \item Más poder, menos sobrantes
\end{itemize}

\subsection{La Clase Disparo - Orquestadora}

La clase Disparo actúa como orquestadora de todas las estrategias. Mantiene una lista de qué estrategias están disponibles y cuál es la actualmente seleccionada. Cuando se quiere disparar, consulta la estrategia actual, le pide que cree un disparo, y lo añade a la lista de disparos activos. También maneja el cambio entre tipos de disparo, asegurándose de seleccionar solo aquellos que tienen munición disponible.

\subsection{Diferencias entre Naves}

Cada nave tiene un arsenal diferente. La Nave1 puede usar píxeles y flechas. La Nave2 puede usar píxeles y rombos. La Nave3, siendo la más potente, tiene acceso a los tres tipos de disparo. Esto hace que la elección de nave no sea solo estética, sino que afecte estratégicamente al juego.

\subsection{¿Qué Ocurre cuando Disparas?}

Cuando el jugador presiona el botón de disparo, el sistema chequea cuál es la estrategia actualmente seleccionada. Si tiene munición disponible, le pide a esa estrategia que cree el disparo con su forma característica. El nuevo disparo se añade a la lista de disparos activos. Cada frame, todos los disparos se mueven hacia arriba, y si salen de la pantalla, simplemente desaparecen. Es un proceso simple y limpio gracias a que cada estrategia sabe cómo se ve su disparo.

\subsection{Ventajas del Patron Strategy}

\begin{enumerate}
    \item \textbf{Intercambiabilidad en tiempo de ejecucion:} El jugador puede cambiar entre disparos dinamicamente.
    
    \item \textbf{Encapsulacion de comportamientos:} Cada estrategia gestiona su municion y crea su forma especifica.
    
    \item \textbf{Eliminacion de condicionales complejos:} Sin Strategy habría complejos \texttt{if-else} esparcidos por el codigo.
    
    \item \textbf{Extensibilidad:} Nueva estrategia (ej. \texttt{DisparoLaser}) solo requiere implementar la interfaz.
    
    \item \textbf{Testabilidad:} Cada estrategia se prueba independientemente.
\end{enumerate}

\newpage
\section{Composite: Armando Todo junto}

\subsection{La Idea}

Una nave está hecha de píxeles. Un disparo complejo también está hecho de píxeles. Composite permite tratar ambos de forma uniforme: un píxel individual se comporta como cualquier grupo de píxeles. Cuando le dices a una nave que se mueva, automáticamente todos sus píxeles se mueven.

Lo práctico de esto:
\begin{itemize}
    \item Naves y disparos complejos se construyen de forma flexible
    \item El movimiento es automático: mueves el contenedor y el contenido se mueve junto
    \item Detección de colisiones: simplemente comparas píxeles
\end{itemize}

\subsection{Cómo está implementado}

Todas las estructuras (píxeles, naves, disparos) implementan la interfaz \texttt{Component}. Un \texttt{Pixel} es la unidad más básica. Un \texttt{Composite} es un contenedor que puede tener píxeles u otros Composites dentro, permitiendo una jerarquía arbitraria.

\subsection{La Interfaz Component}

Todas las estructuras, sean píxeles individuales o composites complejos, implementan la interfaz Component. Esta interfaz establece operaciones comunes: moverse, obtener la posición de referencia, verificar si el componente sigue activo, y notificar eventos importantes al tablero. De esta manera, sin importar si trabajas con un píxel individual o con una estructura compleja, los métodos funcionan igual.

\subsection{Pixel - La Unidad Básica}

Un píxel es simplemente una celda en el tablero. Tiene una posición (X, Y). Puede estar activo o inactivo. Lo interesante es que existen dos tipos: píxeles que forman parte de una nave, y píxeles que forman parte de un disparo. Se comportan diferente. Los de disparo notifican cada movimiento para detectar colisiones, mientras que los de nave solo se mueven. Los de disparo además se desactivan automáticamente si salen de la pantalla.

\subsection{Composite - El Contenedor}

Composite es simplemente un contenedor. Dentro puede tener píxeles individuales, u otros Composites. Eso permite jerarquías arbitrarias. Lo poderoso es que Composite también implementa Component, así que se comporta igual que un píxel. Cuando le dices que se mueva, él se mueve y automáticamente mueve todo lo que tiene dentro.

Existen dos modos de funcionamiento. En modo "nave", valida que todos los píxeles quepan dentro del tablero antes de permitir el movimiento. En modo "disparo", se mueve libremente pero sus píxeles se desactivan si salen del límite superior.

Cuando una nave se mueve, notifica el movimiento de todos sus píxeles en bloque al sistema. Esto es más eficiente que notificar cada píxel por separado.

\subsection{Construyendo Formas}

Cada nave tiene una geometría diferente. Nave1 es como una T invertida (10 píxeles). Nave2 es más rectangular (8 píxeles). Nave3 es compacta, solo 4 píxeles. Simplemente se llena el Composite con los píxeles en las posiciones necesarias.

Con los disparos es lo mismo. Pixel simple = un punto. Flecha = 3 puntos en forma de V. Rombo = 5 puntos en forma de rombo. Cada estrategia de disparo crea su forma adicionando píxeles al Composite.

\subsection{Moviendo una Nave}

Cuando el jugador presiona una tecla, la nave intenta moverse. Primero valida que los píxeles quepan. Si alguien se saldría del borde, no se mueve nada. Si está bien, mueve todos los píxeles automáticamente. Notifica al sistema los cambios para que pueda chequear colisiones. Todo ocurre de forma transparente.

\newpage
\section{Cómo Todo Funciona Junto}

\subsection{El Flujo Completo}

Al iniciar el juego, Factory crea la nave elegida. Esa nave tiene su estructura de píxeles (Composite) y sus disparos disponibles (Strategies). Durante el juego, todo se automatiza: movimientos de naves, disparos volando, colisiones. Los tres patrones trabajan juntos de forma transparente, cada uno resolviendo exactamente su problema.

\subsection{SOLID sin Ser Complicado}

La buena noticia es que al usar estos patrones correctamente, naturalmente terminas teniendo:
\begin{itemize}
    \item Responsabilidades claras: cada clase hace una cosa
    \item Fácil de extender sin romper lo existente
    \item Componentes intercambiables
    \item Dependencias limpias
    \item Código testable
\end{itemize}

\subsection{Beneficios Prácticos}

\begin{itemize}
    \item Agregar una nueva nave: copiar, renombrar, cambiar formas
    \item Agregar un nuevo disparo: implementar Strategy, listo
    \item Cambiar velocidades o posiciones: edita Factory
    \item Detectar colisiones: simple, solo comparas píxeles
    \item Futuras características: la arquitectura aguanta fácilmente
\end{itemize}

Es un diseño que deja espacio para crecer sin colapsar bajo su propio peso.

\newpage
\section{Conclusión}

Los tres patrones utilizados en Space Invaders resuelven problemas de diseño concretos y reales. Factory centraliza la creación, Strategy permite cambiar comportamientos dinámicamente, y Composite proporciona una estructura flexible y uniforme. Juntos, crean una arquitectura que es fácil de mantener, entender y extender. Cada patrón aporta su parte, y el resultado es un código limpio, organizado y preparado para el futuro.

\end{document}
